\begin{lemma}
\label{editorial-source-lemma-map-into-tensor-algebra-open}
Let $k$ be a field.
Let $R$ and $S$ be $k$-algebras.
The map $\Spec(S \otimes_k R) \to \Spec(R)$
is open.
\end{lemma}

\begin{proof}
Retain the original $k$-algebras $S,R$ and write the
specified tensor as a finite sum
$$
f=\sum_{i=1}^d s_i\otimes r_i\in S\otimes_kR.
$$
Set $S'=k[s_1,\ldots,s_d]\subseteq S$, keeping every
original coefficient $s_i$. Tensoring $S'\hookrightarrow S$
with $R$ is injective because $R$ is a $k$-vector space,
so the same displayed tensor is an element of
$S'\otimes_kR$ with its specified image $f$.
The $k$-algebra $S'$ is finitely presented: it is
a quotient of $k[x_1,\ldots,x_d]$, and this polynomial
ring is Noetherian by Lemma \ref{editorial-source-lemma-Noetherian-permanence}.
Hence the kernel of the original map $x_i\mapsto s_i$
has finitely many generators $h_1,\ldots,h_m$.
This also follows from
Lemma \ref{editorial-source-lemma-Noetherian-finite-type-is-finite-presentation}.

The base change retains the presentation
$$
S'\otimes_kR\simeq R[x_1,\ldots,x_d]/(h_1^R,\ldots,h_m^R),
$$
where every coefficient of $h_j^R$ is the image of
the corresponding coefficient of $h_j$ in $R$.
The map sends $s_i\otimes1$ to the class of $x_i$
and $1\otimes r$ to $r$, and its inverse sends the
original polynomial generators and coefficients back
to these tensors. The relations and balancing prove
both maps are well defined and inverse.
The algebra is flat over $R$. To see this directly,
choose a $k$-basis $(b_\lambda)$ of $S'$.
The map
$$
\bigoplus_\lambda R\longrightarrow S'\otimes_kR,
\qquad (a_\lambda)_\lambda\longmapsto
\sum_\lambda b_\lambda\otimes a_\lambda
$$
is an $R$-module isomorphism; its inverse is the
finite basis expansion in the first tensor factor.
Tensoring a free module preserves injections by
the corresponding direct-sum maps. Thus
Proposition \ref{editorial-source-proposition-fppf-open} applies to
$R\to S'\otimes_kR$ and makes the image of the
original principal open $D(f)$ in this spectrum
open. Lemma \ref{editorial-source-lemma-same-image}, with the exact
same $f$, identifies it with the image of $D(f)$
in $\Spec(S\otimes_kR)$. Principal opens form a
basis and images commute with unions, so the
original spectrum map is open, as claimed.

This also proves universal openness of the
original map. For an arbitrary original $R$-algebra
$A$, its base change is identified with
$\Spec(S\otimes_kA)\to\Spec(A)$ through the map
$$
(S\otimes_kR)\otimes_RA\longrightarrow S\otimes_kA,
\qquad (s\otimes r)\otimes a\longmapsto s\otimes\overline r a,
$$
whose inverse is $s\otimes a\mapsto(s\otimes1)\otimes a$.
Both maps respect the original algebra products
and balancing, and their composites fix all
elementary tensors. The result already proved
applies to this $k$-algebra $A$, giving openness
after every base change, without finite generation
of the full algebra $S$.

\medskip\noindent
One can make the principal-open image more explicit
without confusing nonzero and nonnilpotent elements.
For the fixed original list $s_1,\ldots,s_d$, there
exists an integer $N\geq1$, independent of $R$
and $r_1,\ldots,r_d$, such that the image of $D(f)$
is the union of the coefficient opens of the full
power $f^N$ in a $k$-basis of $S'$.
Here is a proof with all coefficients retained.

Put $A_0=k[T_1,\ldots,T_d]$ and
$F=\sum_{i=1}^d s_i\otimes T_i\in S'\otimes_kA_0$.
Choose a $k$-basis $(b_\lambda)_{\lambda\in\Lambda}$
of $S'$. For each ordered tuple $(i_1,\ldots,i_n)$,
write the complete finite expansion
$$
s_{i_1}\cdots s_{i_n}
=\sum_{\lambda\in\Lambda}
\mu_{i_1,\ldots,i_n,\lambda}b_\lambda,
\qquad \mu_{i_1,\ldots,i_n,\lambda}\in k.
$$
Then, with all ordered tuples and their repeated
terms retained, the coefficients of $F^n$ are
$$
\begin{split}
F^n&=\sum_{\lambda\in\Lambda}b_\lambda\otimes C_{n,\lambda}(T),\\
C_{n,\lambda}(T)
&=\sum_{i_1=1}^d\cdots\sum_{i_n=1}^d
\mu_{i_1,\ldots,i_n,\lambda}T_{i_1}\cdots T_{i_n}.
\end{split}
$$
For fixed $n$ only finitely many of these polynomials
are nonzero. Put $J_n=(C_{n,\lambda}:\lambda\in\Lambda)
\subseteq A_0$ and $Z_n=V(J_n)$. Extension of the
basis to every residue field shows that
$\mathfrak p\in Z_n$ exactly when the original
$F$ has $n$-th power zero in
$S'\otimes_k\kappa(\mathfrak p)$. Hence
$Z_n\subseteq Z_{n+1}$. The fibre criterion and
the tensor-localization maps of the preceding
lemma give
$$
\Spec(A_0)\setminus
\operatorname{image}\bigl(D(F)\to\Spec(A_0)\bigr)
=\bigcup_{n\geq1}Z_n.
$$
Call this complement $Z$. The proven openness
result, applied to the original algebra $S'$
and base $A_0$, says that $Z$ is closed.

The ring $A_0$ is Noetherian, so its closed set
$Z$ has finitely many irreducible components
$V(\mathfrak p_1),\ldots,V(\mathfrak p_e)$ with
their indicated generic points. This follows
by applying Lemma \ref{editorial-source-lemma-Noetherian-irreducible-components}
to the quotient defining $Z$. Each generic point
lies in some $Z_{n_j}$, because their union is
$Z$. Since $Z_{n_j}$ is closed, it contains
the full component $V(\mathfrak p_j)$.
Thus for $N=\max(1,n_1,\ldots,n_e)$,
$$
Z=Z_N.
$$
If $Z$ is empty, use $N=1$: every $Z_n$ is
already empty. This is a finite-component argument;
it does not infer stabilization of increasing
closed sets merely from Noetherianity.

For any original $R$ and $r_i$, use the original
map $A_0\to R$, $T_i\mapsto r_i$.
At a prime $\mathfrak p\subset R$, its contraction
$\mathfrak p_0\subset A_0$ induces an injective
field map $\kappa(\mathfrak p_0)\to\kappa(\mathfrak p)$.
The map
$S'\otimes_k\kappa(\mathfrak p_0)\to
S'\otimes_k\kappa(\mathfrak p)$ is injective by
extension of vector spaces, and carries the full
original $F$ to $f$. Therefore nilpotence is
equivalent in these two fibres, and by $Z=Z_N$
it is equivalent to the vanishing of the $N$-th
power. Its coefficients are exactly
$C_{N,\lambda}(r_1,\ldots,r_d)$, with the same
ordered-tuple formula. It follows that
$$
\operatorname{image}\bigl(D(f)\subseteq
\Spec(S\otimes_kR)\to\Spec(R)\bigr)
=\bigcup_{\lambda\in\Lambda}
D\bigl(C_{N,\lambda}(r_1,\ldots,r_d)\bigr).
$$
Only finitely many original coefficient polynomials
contribute, giving a finite union of basic opens
and hence a quasi-compact open image. The equality
for $S$ uses the same-image lemma after the equality
for its specified subalgebra $S'$. None of these
comparisons replaces $f$ by its linear nonzero
coefficient test: the required test uses the
proved power $N$. For instance in
$S'=k[\epsilon]/(\epsilon^2)$ with $f=\epsilon\otimes1$,
the first power is nonzero and has a unit coefficient,
while the actual principal open is empty; the
second power vanishes, giving the correct empty
union. The case $d=0$ has $F=f=0$ and can use
$N=1$, and zero tensor algebras give empty images
under the same formulas.
\end{proof}